% !TeX spellcheck = ru_RU
% !TEX root = vkr.tex

\section{Постановка задачи}
\label{sec:task}

\subsection{Цель и задачи работы}

Работа выполняется в интересах лаборатории Гербарий ГБС РАН, сотрудники которой проводят морфометрический анализ листьев мохообразных (в первую очередь мхов рода \textit{Sphagnum}).

Рабочий процесс биолога включает два принципиально разных этапа.
Первый~--- \textit{отбор образцов}: на каждом изображении присутствуют несколько листьев,
и исследователь выбирает те из них, которые пригодны для измерений.
«Хороший» лист~--- целый, с неповреждёнными краями, типичной для вида формы,
расположен изолированно и не перекрывается с другими объектами.
Листья с рваными или нечёткими краями, артефактами препарирования,
атипичной морфологией или сильным перекрытием исключаются.
Критерии отбора субъективны, определяются целями конкретного исследования
и требуют экспертной оценки биолога; формализовать их в виде алгоритма невозможно.
Второй этап~--- \textit{сегментация и измерение} отобранных листьев~---
представляет собой технически однообразную ручную работу, которая при сериях из десятков
и сотен образцов занимает значительное время и вносит субъективную погрешность.

\textbf{Цель работы}~--- разработать настольное приложение, автоматизирующее цикл
«загрузка изображений → сегментация → ручная коррекция → параметрические измерения → экспорт»,
пригодное для использования без специальных вычислительных ресурсов на рабочих станциях биологической лаборатории.

Для достижения цели поставлены следующие \textbf{задачи}:
\begin{enumerate}
    \item Проанализировать существующие программные решения для сегментации и морфометрии биологических объектов и обосновать выбор модели сегментации.
    \item Сформулировать требования к системе совместно с заказчиком~--- сотрудниками лаборатории Гербарий ГБС РАН.
    \item Спроектировать архитектуру приложения и реализовать его.
    \item Реализовать алгоритм параметрических измерений (длина, поперечные профили ширины).
    \item Сравнить результаты автоматической сегментации с ручной разметкой биологов и провести апробацию системы.
\end{enumerate}

\textbf{Практическая значимость} работы состоит в сокращении времени, затрачиваемого на получение морфометрических данных, и в повышении воспроизводимости измерений за счёт стандартизованного алгоритма обработки.

